% ══════════════════════════════════════════════════════════════════════════════
% CHAPTER 1 — INTRODUCTION
% ══════════════════════════════════════════════════════════════════════════════

\section{Motivation}
\label{sec:intro_motivation}

Central bank communication has become an important instrument of
monetary policy. Modern central banks influence expectations not
only through policy-rate decisions, but also through the language
used in post-meeting statements, press conferences, minutes, and
other public communication. Understanding \emph{what} central banks
communicate and \emph{how} financial markets interpret this
communication is therefore central to monetary economics and
financial markets research \citep{blinder2008}.

This role became particularly visible during the era of
unconventional monetary policy. At the zero lower bound, forward
guidance allowed central banks to influence expectations about the
future path of policy when conventional rate cuts were constrained
\citep{campbell2012, swanson2021}. More generally,
\citet{gurkaynak2005} show that FOMC announcements cannot be
summarised by the current rate surprise alone. Instead, a separate
path factor, closely related to communication about future policy,
explains a substantial part of asset-price reactions, especially at
longer Treasury maturities. This motivates the empirical design of
controlling for the target-rate surprise while testing whether
textual tone explains additional market reactions.

Measuring central bank communication, however, remains
methodologically challenging. Dictionary-based approaches are
transparent and easy to interpret, but rely on predefined word lists
and can miss context. For example, references to asset purchases are
not inherently dovish or hawkish: ``increasing asset purchases'' and
``decreasing asset purchases'' imply opposite policy signals,
although a simple word-count approach may treat them similarly. This
relates to the broader concern that dictionaries can misclassify
domain-specific language when word meanings depend on context
\citep{loughran2011}.

Transformer-based models such as RoBERTa provide a contextual
alternative. \citet{shah2023} introduce the \emph{Trillion Dollar
Words} dataset and show that fine-tuned language models can classify
FOMC sentences as hawkish, dovish, or neutral. These models offer a
useful benchmark, although their predictions are less transparent
than dictionary counts and depend on the quality of the fine-tuning
data.

This thesis compares three complementary approaches on the same
corpus of central bank texts: a hand-crafted dictionary method, a
class-balanced Random Forest classifier, and a fine-tuned RoBERTa
sentence-level classifier. This comparison allows the thesis to
assess how transparent rule-based methods, supervised machine
learning, and contextual transformer models differ in measuring the
tone of central bank communication.

\section{Research Questions}
\label{sec:intro_rqs}

This thesis investigates two primary research questions:

\begin{description}[leftmargin=1.5cm, style=nextline]
  \item[RQ1] Does the hawkish--dovish tone of FOMC and ECB
    communications explain same-day financial market reactions (maybe change to announcement-window financial market reactions chatgpt ?) in
    equity markets, foreign exchange markets, government bond yields,
    sovereign spreads, and inflation breakevens, after controlling for
    the monetary policy rate surprise?

  \item[RQ2] Can central bank language be used to forecast near-term
    monetary policy outcomes and yield changes in an expanding-window
    setting, beyond simple policy inertia?
\end{description}

RQ1 is examined through event-study regressions that relate same-day
financial market reactions (mb change to RQ1 is examined through event-study regressions that relate
announcement-window financial market reactions chatgpt? ) to document-level hawkish--dovish tone
scores, while controlling for the monetary policy rate surprise. RQ2 is
examined through expanding-window forecasting exercises, in which
dictionary-based forward-looking tone measures, Random Forest models,
and RoBERTa-based tone scores are used to predict future rate decisions
and yield changes.

In addition to the dictionary approach, the thesis develops two
data-driven extensions. The Random Forest produces an alternative
continuous tone measure based on out-of-bag class probabilities, while
the RoBERTa classifier aggregates sentence-level hawkish, dovish, and
neutral predictions into a document-level score using the same
alpha-smoothed net-score formula as the dictionary. Horse-race
regressions including all three measures simultaneously test whether
the different methods capture independent market-relevant information.

\section{Contribution}
\label{sec:intro_contribution}

This thesis contributes to the literature in six ways.

First, it provides a \textbf{dual-bank, long-sample comparison} using a
consistent methodology for both the FOMC and the ECB. FOMC
post-meeting statements span February~1994 to March~2026, while ECB
communications span January~1999 to March~2026. The combined sample
covers several monetary policy regimes, including the Great Moderation,
the Global Financial Crisis, the zero lower bound, and the post-COVID
tightening cycle. This long horizon makes it possible to examine whether
the relationship between central bank communication and financial market
reactions is robust across different policy environments.

Second, it introduces a \textbf{structured dictionary decomposition} of
central bank tone. Rather than relying only on a single net
hawkish--dovish score, the dictionary separates communication into
policy action, inflation, employment, and residual risk language. These
sub-scores enter the market regressions under three nested
specifications: a collapsed net score, a policy/economic split, and the
full four-component decomposition. This allows the analysis to examine
which dimensions of communication are most closely associated with
different financial market reactions.

Third, it applies \textbf{three complementary scoring methods to the
same corpus} of 884 central bank documents: a self-crafted dictionary
method, a Random Forest classifier using Lasso-based feature selection,
and a fine-tuned RoBERTa sentence-level classifier \citep{shah2023}.
Horse-race regressions, in which all three scores enter simultaneously,
test whether the dictionary, Random Forest, and transformer-based
measures capture independent market-relevant information.

Fourth, it introduces a \textbf{forward-looking versus backward-looking
sentence decomposition}. Sentences are classified using linguistic
markers related to futurity, commitment, projections, and
conditionality. This decomposition is motivated by the distinction
between current policy surprises and future path-of-policy information
in \citet{gurkaynak2005}, but implements it directly at the textual
level. The resulting forward-looking dictionary and RoBERTa-based scores
allow the thesis to test whether forward-guidance-related language has a
distinct association with market reactions and forecasting performance.

Fifth, the Random Forest pipeline incorporates several
\textbf{methodological refinements} for classifying central bank
documents. It uses bigram TF-IDF features to capture phrase-level stance
signals, Lasso-based pre-filtering to reduce the dimensionality of the
high-dimensional text space, and class-balancing to address the
dominance of unchanged policy decisions. This design allows the model to
learn nonlinear textual patterns while limiting overfitting.

Sixth, the forecasting analysis in RQ2 employs a \textbf{strictly
out-of-sample expanding-window design}. The models forecast four
targets: the rate decision at the next meeting, the rate decision two
meetings ahead, the two-year yield change at the next meeting, and the
two-year yield change two meetings ahead. At each meeting~$t$, feature
selection and model estimation are performed using only documents
available up to $t-1$. This design prevents look-ahead bias and provides
a real-time test of whether central bank language contains forecasting
information beyond simple policy inertia.
Reviewed AP


\section{Scope and Delimitation}
\label{sec:intro_scope}
The empirical analysis focuses on the FOMC and the ECB. This scope is
chosen because both institutions are covered by publicly available
high-frequency monetary policy event databases: the USMPD for the
United States and the EA-MPD for the euro area. These datasets measure
changes in financial variables around narrow monetary policy event
windows, including statement windows, press conference windows, and
broader monetary policy announcement windows. This makes them suitable
for studying whether central bank communication is associated with
immediate financial market reactions.

The empirical framework consists of separate event-level regressions for
each central bank and communication corpus. In these regressions, the
dependent variables are high-frequency market reactions around the
relevant event window, while the main explanatory variables are
document-level hawkish--dovish tone measures. The regressions control
for the monetary policy rate surprise, allowing the analysis to test
whether communication tone explains market movements beyond the
mechanical surprise component of the policy decision. Estimating the
models separately for the FOMC and the ECB preserves the
institution-specific interpretation of the communication effect, since
the two central banks differ in institutional structure, communication
format, and event-window definitions.


\section{Thesis Outline}
\label{sec:intro_outline}
The remainder of this thesis is organised as follows.
\Cref{ch:literature} reviews the literature on central bank
communication, text-based analysis in finance, machine learning
methods, and event-study identification.
\Cref{ch:institutional} provides institutional background on the
FOMC and the ECB and describes the main monetary policy eras covered
by the sample. \Cref{ch:data} presents the textual corpus, financial
market data, and rate decision variables. \Cref{ch:methodology}
introduces the dictionary-based tone measures, the Random Forest and
RoBERTa extensions, and the empirical specifications.
\Cref{ch:results} reports the main findings on tone measurement,
market reactions, forward-looking communication, and forecasting
performance. \Cref{ch:discussion} discusses the results, limitations,
and comparison between the FOMC and ECB. \Cref{ch:conclusion}
concludes.